\subsection{Boundary comparison from actual local maxima}\label{subsec:fc-line-boundary-comparison}

\begin{lemma}[Compact connected boundary comparison]\label{lem:fc-line-boundary-comparison}
\lean{EnvelopeProperty.Berkovich.RelativeSpectrum.FCStandardAffinoid.le_zero_of_boundary_of_local_max_plateau}
\leanok
\uses{def:fc-standard-affinoid, lem:fc-line-topology}
Let $U$ be an actual standard affinoid and let $f:U\to[-\infty,+\infty]$
be continuous. Suppose $f\le0$ on the actual Gauss boundary. Suppose also
that every positive interior local maximum of the ambient function $f$
has an ambient neighbourhood on which $f$ is constant. Then $f\le0$ on $U$.
\end{lemma}

\begin{proof}
If a positive value occurs, compactness supplies a positive global maximum.
Its locus is nonempty and closed by continuity, and avoids the boundary.
Every point of this locus is therefore an interior local maximum. The
assumed local constancy makes the locus open relative to $U$. Connectedness
then makes it all of $U$, contradicting the boundary inequality at the outer
Gauss point. This is a general topological lemma; subsequent applications
derive its local hypothesis from the actual rational-unit maximum theorem.
\end{proof}
